% Chapter 6: Headlamps
\chapter{前照灯安装规定对比}

\section{远光灯}

\subsection{安装要求（基于实际法规文本）}

\begin{table}[H]
\centering
\caption{远光灯安装要求对比（基于实际法规文本）}
\small
\begin{tabularx}{\textwidth}{|l|X|X|X|}
\hline
\textbf{项目} & \textbf{\versionAShort (2019 5.1)} & \textbf{\versionBShort (草案 5.1)} & \textbf{\versionCShort (UN R48-09 6.1)} \\
\hline
配备 & 汽车应配备，挂车不应使用 & 同左 & 同左 \\
\hline
数量 & 2 只或 4 只（N3 类可 + 2 只）& 同左 & 同左 \\
\hline
颜色 & 白色 & 同左 & White \\
\hline
安装高度 & 无特殊要求 & 不应高于前挡风玻璃下边缘 & 无特殊要求 \\
\hline
横向间距 & 无特殊要求 & 无特殊要求 & 无特殊要求 \\
\hline
纵向 & 装在车前 & 同左 & 装在车前 \\
\hline
ADB 集成 & 5.1.7.1 自动控制 & 同左 & 5.7 自动控制 \\
\hline
\end{tabularx}
\end{table}

\subsection{光度要求}

\begin{table}[H]
\centering
\caption{远光灯光度要求对比（基于实际法规文本，单位：cd）}
\begin{tabular}{|l|c|c|c|}
\hline
\textbf{测试点} & \textbf{\versionAShort (5.1.9.1)} & \textbf{\versionBShort} & \textbf{\versionCShort} \\
\hline
最大发光强度之和 & $\le 430,000$ & 同左 & 同左 \\
\hline
HV 点（参考轴） & $\ge 18000$, $\le 300000$ & 同左 & $\ge 18000$, $\le 300000$ \\
\hline
HV+0.5L, HV+0.5R & $\ge 12000$ & 同左 & $\ge 12000$ \\
\hline
HV+1L, HV+1R & $\ge 6000$ & 同左 & $\ge 6000$ \\
\hline
HV+1.5L, HV+1.5R & $\ge 2000$ & 同左 & $\ge 2000$ \\
\hline
HV+2L, HV+2R & $\ge 1000$ & 同左 & $\ge 1000$ \\
\hline
HV+0.5U, HV+0.5D & $\ge 6000$ & 同左 & $\ge 6000$ \\
\hline
\end{tabular}
\end{table}

\subsection{远光灯自动控制要求（5.1.9.3）}

\begin{table}[H]
\centering
\caption{远光灯自动控制要求对比（基于实际法规文本）}
\begin{tabular}{|l|c|c|c|}
\hline
\textbf{项目} & \textbf{\versionAShort} & \textbf{\versionBShort} & \textbf{\versionCShort} \\
\hline
传感器水平角度 & $\pm 15°$ & 同左 & 同左 \\
\hline
迎面来车探测距离 & $\ge 400$ m & $\ge 400$ m & 同左 \\
\hline
前方车辆探测距离 & $\ge 100$ m & $\ge 100$ m & 同左 \\
\hline
电动自行车探测距离 & 未提及 & $\ge 75$ m（草案新增） & --- \\
\hline
环境照度自动关闭阈值 & $> 7000$ lx & $> 7000$ lx & 同左 \\
\hline
\end{tabular}
\end{table}

\section{近光灯}

\subsection{安装要求（基于实际法规文本）}

\begin{table}[H]
\centering
\caption{近光灯安装要求对比（基于实际法规文本）}
\small
\begin{tabularx}{\textwidth}{|l|X|X|X|}
\hline
\textbf{项目} & \textbf{\versionAShort (2019 5.2)} & \textbf{\versionBShort (草案 5.2)} & \textbf{\versionCShort (UN R48-09 6.2)} \\
\hline
配备 & 汽车应配备，挂车不应使用 & 同左 & 同左 \\
\hline
数量 & 2 只 & 同左 & 2 只 \\
\hline
横向：离外缘端面距离 & $\le 400$ mm & 同左 & $\le 400$ mm \\
\hline
横向：两灯间距 & $\ge 600$ mm（M1/N1 无；车宽 $< 1300$ 时 $\ge 400$ mm）& 同左 & 同左 \\
\hline
高度 & $500$--$1200$ mm（N3G 类可增至 $1500$ mm）& 同左 & 同左 \\
\hline
纵向 & 装在车前 & 同左 & 装在车前 \\
\hline
可见度 $\alpha$ 角 & 上 $15°$ 下 $10°$ & 同左 & 上 $15°$ 下 $10°$ \\
\hline
可见度 $\beta$ 角 & 外 $45°$ 内 $10°$ & 同左 & 外 $45°$ 内 $10°$ \\
\hline
明暗截止线 & 必须清晰 & 同左 & 必须清晰 \\
\hline
\end{tabularx}
\end{table}

\subsection{光度要求}

\begin{table}[H]
\centering
\caption{近光灯光度要求对比（单位：lx）}
\begin{tabular}{|l|c|c|c|}
\hline
\textbf{测试点} & \textbf{\versionAShort} & \textbf{\versionBShort} & \textbf{\versionCShort} \\
\hline
HV (B50L) & $\le 350$ lx (防眩目) & 同左 & $\le 350$ lx \\
\hline
75R & $\ge 12000$ cd & 同左 & $\ge 12000$ cd \\
\hline
50R & $\ge 12000$ cd & 同左 & $\ge 12000$ cd \\
\hline
25R & $\ge 6000$ cd & 同左 & $\ge 6000$ cd \\
\hline
25L & $\le 1200$ cd (防眩目) & 同左 & $\le 1200$ cd \\
\hline
50V & $\ge 6000$ cd & 同左 & $\ge 6000$ cd \\
\hline
III 区 (防眩目) & $\le 625$ cd & 同左 & $\le 625$ cd \\
\hline
\end{tabular}
\end{table}

\subsection{明暗截止线具体规定}

\begin{table}[H]
\centering
\caption{明暗截止线规定对比}
\small
\begin{tabularx}{\textwidth}{|l|X|X|X|}
\hline
\textbf{项目} & \textbf{\versionAShort} & \textbf{\versionBShort} & \textbf{\versionCShort} \\
\hline
清晰度 & 必须可清晰分辨 & 强化，具体量化要求 & 必须可清晰分辨 \\
\hline
水平部分 & 必须水平 & 同左 & 必须水平 \\
\hline
弯道倾斜 & 允许"弯道倾斜" 15° & 同左 & 加强弯道定义 \\
\hline
ADB 状态下的变化 & 未规定 & 新增：ADB 模式下的截止线动态调节要求 & ADB 截止线调整 \\
\hline
\end{tabularx}
\end{table}

\section{AFS（自适应前照灯系统）}

\subsection{AFS 等级对比}

\begin{table}[H]
\centering
\caption{AFS 等级划分对比}
\begin{tabular}{|l|c|c|c|}
\hline
\textbf{等级} & \textbf{\versionAShort} & \textbf{\versionBShort} & \textbf{\versionCShort} \\
\hline
C 等级（基本近光） & $\checkmark$ & $\checkmark$ & $\checkmark$ \\
\hline
V 等级（市区近光） & $\checkmark$ & $\checkmark$ & $\checkmark$ \\
\hline
E 等级（高速近光） & $\checkmark$ & $\checkmark$ & $\checkmark$ \\
\hline
W 等级（恶劣天气近光） & $\checkmark$ & $\checkmark$ & $\checkmark$ \\
\hline
P 等级（前雾灯） & $\checkmark$ & $\checkmark$ & $\checkmark$ \\
\hline
弯道照明等级 & $\checkmark$ & 强化 & $\checkmark$ \\
\hline
ADB 等级（远光自适应） & 未规定 & 新增 & $\checkmark$ \\
\hline
\end{tabular}
\end{table}

\subsection{AFS 切换要求}

\begin{table}[H]
\centering
\caption{AFS 切换要求对比}
\small
\begin{tabularx}{\textwidth}{|l|X|X|X|}
\hline
\textbf{项目} & \textbf{\versionAShort (2019 5.22)} & \textbf{\versionBShort (草案 5.3)} & \textbf{\versionCShort (UN R48-09 6.22)} \\
\hline
等级切换触发条件 & 车速/转向角/环境感知 & 强化，新增"环境感知"参数 & Speed/steering \\
\hline
切换时间 & 一般规定 & $\le 1.5$ 秒（强化） & 一般规定 \\
\hline
指示 & 黄色 AFS 指示灯 & 同左，新增"故障降级指示" & Yellow tell-tale \\
\hline
故障检测 & 一般规定 & 强化，新增"自诊断"要求 & Self-diagnostic \\
\hline
降级模式 & 一般规定 & 强化，明确"近光基本等级" & Degraded mode \\
\hline
\end{tabularx}
\end{table}

\subsection{AFS 照明单元位置要求（基于实际条款）}

基于 GB 4785-2019 第 5.22.4.1 条与草案版 5.3.4.1 条的相同规定，AFS 照明单元之间的几何关系如下：

\begin{table}[H]
\centering
\caption{AFS 照明单元位置参数对比}
\begin{tabular}{|l|c|c|}
\hline
\textbf{参数} & \textbf{实际值（基于图 3）} & \textbf{说明} \\
\hline
水平方向 A & $\le 400$ mm & 单元视表面外缘到车辆外缘端面的距离 \\
\hline
水平方向 B & $\ge 600$ mm（M1/N1 无；车宽 $< 1300$ 时 $\ge 400$ mm）& 两视表面内边缘间的距离 \\
\hline
水平方向 C & $\le 200$ mm & 单元之间的最大水平距离 \\
\hline
水平方向 E & $\le 140$ mm & 附加照明单元与相邻照明单元之间的最大水平距离 \\
\hline
垂直方向 D & $\le 400$ mm & 与上/下相邻照明单元之间的垂直距离 \\
\hline
垂直方向 F & $\ge 250$ mm & 附加照明单元安装高度下限 \\
\hline
垂直方向 G & $\le 1200$ mm & 符合近光灯最大离地高度要求 \\
\hline
对称安装误差 & $\le 50$ mm & 对称安装的照明单元视表面重心到车辆纵向对称平面的距离误差 \\
\hline
\end{tabular}
\end{table}

\subsection{AFS 调光装置（基于实际条款 5.22.6 / 6.22.6）}

\begin{table}[H]
\centering
\caption{AFS 调光装置要求对比}
\begin{tabular}{|l|c|c|c|}
\hline
\textbf{项目} & \textbf{\versionAShort} & \textbf{\versionBShort} & \textbf{\versionCShort} \\
\hline
基础近光明暗截止线初始下倾度 & $0.1\%$ 准确度标注 & 同左 & $0.1\%$ 准确度标注 \\
\hline
自动调光装置 & 强制 & 同左 & 强制 \\
\hline
弯道照明 & 允许 & 同左 & 允许 \\
\hline
弯道照明轨迹半径 & $\le 500$ m & 同左 & $\le 500$ m \\
\hline
\end{tabular}
\end{table}

\section{ADB（自适应远光灯）—— 草案版新增重点}

\subsection{基本定义}

\textbf{ADB（Adaptive Driving Beam）} 是指能根据前方车辆、交通环境、道路条件自动调整远光光型，以避免对对向和前方车辆造成眩目的前照灯系统。

\subsection{技术要求}

\begin{table}[H]
\centering
\caption{ADB 技术要求对比}
\small
\begin{tabularx}{\textwidth}{|l|X|X|X|}
\hline
\textbf{项目} & \textbf{\versionAShort} & \textbf{\versionBShort} & \textbf{\versionCShort} \\
\hline
适用车辆         & --- & M、N 类车辆 & M, N classes \\
\hline
系统构成           & --- & 摄像头/雷达、控制单元、可控 LED 阵列 & Camera/radar + ECU + LEDs \\
\hline
识别距离          & --- & $\ge 200$ m（前车）  & $\ge 200$ m \\
\hline
切换时间          & --- & $\le 200$ ms          & $\le 200$ ms \\
\hline
眩目控制          & --- & 必须保持非眩目状态  & Non-glare guaranteed \\
\hline
像素分辨率       & --- & $\ge 10000$ 像素（推荐）  & Recommended \\
\hline
失效安全          & --- & 必须进入"基本远光"或"近光"状态 & Fail-safe to DBL or PBL \\
\hline
\end{tabularx}
\end{table}

\subsection{测试方法}

草案版附录 D 中详细规定了 ADB 的测试方法：

\begin{enumerate}[leftmargin=2em]
    \item \textbf{实验室测试}：暗室内模拟各种交通场景，测试 ADB 的响应；
    \item \textbf{路面测试}：在实际道路上测试 ADB 的实际表现；
    \item \textbf{功能安全测试}：测试 ADB 在各种故障情况下的降级模式；
    \item \textbf{EMC 测试}：测试 ADB 的电磁兼容性；
    \item \textbf{环境测试}：测试 ADB 在各种环境条件下的表现。
\end{enumerate}

\subsection{光型分区要求}

ADB 通常将远光光型分为多个独立可控区域。UN R48-09 规定：

\begin{table}[H]
\centering
\caption{ADB 光型分区要求}
\begin{tabular}{|l|c|c|}
\hline
\textbf{分区类型} & \textbf{分区数} & \textbf{功能} \\
\hline
核心区（Central） & 4--8 区 & 主体照明 \\
\hline
边缘区（Peripheral） & 4--8 区 & 道路边缘 \\
\hline
动态屏蔽区（Dynamic shadow） & 任意 & 屏蔽对向/前方车辆 \\
\hline
\end{tabular}
\end{table}

\section{矩阵式 LED 前照灯}

\subsection{基本要求（草案版新增）}

\begin{table}[H]
\centering
\caption{矩阵式 LED 前照灯要求对比}
\small
\begin{tabularx}{\textwidth}{|l|X|X|X|}
\hline
\textbf{项目} & \textbf{\versionAShort} & \textbf{\versionBShort} & \textbf{\versionCShort} \\
\hline
像素数           & --- & $\ge 20$ 像素 & $\ge 20$ 像素 \\
\hline
像素失效控制      & --- & 部分失效不影响功能（$\le 25\%$） & 部分失效可接受 \\
\hline
光型控制          & --- & 像素级独立控制 & Pixel-level control \\
\hline
热管理            & --- & 必须配备散热装置 & Heatsink required \\
\hline
寿命要求          & --- & $\ge 3000$ 小时 & $\ge 3000$ h \\
\hline
测试要求          & --- & 附录 D 新增 & Annex D \\
\hline
\end{tabularx}
\end{table}

\subsection{失效模式}

矩阵式 LED 灯具需要明确失效模式：

\begin{enumerate}[leftmargin=2em]
    \item \textbf{单像素失效}：不影响整体功能，继续工作；
    \item \textbf{多像素失效}：超过 25\% 时必须降级；
    \item \textbf{完全失效}：自动切换到传统远光或近光；
    \item \textbf{指示信号}：驾驶员必须能感知到严重失效。
\end{enumerate}

\section{前照灯清洗器}

\subsection{基本要求对比（基于实际条款）}

\begin{table}[H]
\centering
\caption{前照灯清洗器要求对比（基于实际法规文本）}
\begin{tabular}{|l|c|c|c|}
\hline
\textbf{项目} & \textbf{\versionAShort (5.2.9)} & \textbf{\versionBShort (草案)} & \textbf{\versionCShort (6.2.9)} \\
\hline
强制安装触发阈值 & 单侧光通量 $> 2000$ lm & 同左 & 单侧 $> 2000$ lm \\
\hline
温度范围（可藏装置） & $-30°\text{C}$ 至 $+50°\text{C}$ & 同左 & 同左 \\
\hline
可藏装置达到使用位置时间 & $\le 3$ s & 同左 & $\le 3$ s \\
\hline
\end{tabular}
\end{table}

\section{水平/垂直调节装置}

\subsection{基本要求对比}

\begin{table}[H]
\centering
\caption{水平/垂直调节装置要求对比}
\begin{tabular}{|l|c|c|c|}
\hline
\textbf{项目} & \textbf{\versionAShort} & \textbf{\versionBShort} & \textbf{\versionCShort} \\
\hline
自动调节 & 强制（部分车型） & 强化，新增"ADB 自动调节" & Mandatory (certain classes) \\
\hline
手动调节 & 一般规定 & 同左 & General \\
\hline
调节范围 & 一般规定 & 强化 & 详细规定 \\
\hline
初始调节 & 必须明确 & 同左 & 必须明确 \\
\hline
\end{tabular}
\end{table}

\section{小结}

\begin{enumerate}[leftmargin=2em]
    \item 草案版对前照灯的安装规定整体上向 UN R48-09 进一步对齐；
    \item 草案版新增了 ADB 和矩阵式 LED 的完整规定，填补了现行版的空白；
    \item 草案版对 AFS 的等级切换要求进行了强化；
    \item 草案版对水平/垂直调节装置的要求更加严格，与 UN R48-09 一致；
    \item 这些变化对整车厂和灯具供应商提出了新的技术要求，需要进行相应的产品调整。
\end{enumerate}